problem:  
Let \((X,L)\) be a polarized compact Kähler manifold and \(E^{\mathrm{an}}\) the analytic Mabuchi energy.  For each normal, relatively ample test configuration \((\mathcal{X},\mathcal{L})\) of \((X,L)\), let \(\varphi_t\) denote the associated weak geodesic ray in the space of Kähler potentials.  It is known (Boucksom–Hisamoto–Jonsson) that  
\[
E^{\mathrm{an}}(\varphi_t)=tE^{\mathrm{NA}}(\mathcal{X},\mathcal{L})+o(t)\,.
\]
A natural conjectural refinement would be the existence of a **universal second‑order expansion**  
\[
E^{\mathrm{an}}(\varphi_t)=tE^{\mathrm{NA}}(\mathcal{X},\mathcal{L})+C(\mathcal{X},\mathcal{L})+o(1)
\tag{★}
\]
where \(C(\mathcal{X},\mathcal{L})\in\mathbb{R}\) depends only on the equivalence class of the test configuration.  

**Question.**  
1. Does the expansion (★) hold for all normal test configurations of a polarized Kähler manifold?  
2. If so, is the constant term \(C(\mathcal{X},\mathcal{L})\) invariant under base change and equivariant isomorphism of test configurations?  
3. Can \(C(\mathcal{X},\mathcal{L})\) be expressed in purely algebraic terms—for instance via an intersection number on the total space of \(\mathcal{X}\) involving the relative canonical divisor—or interpreted as a non‑Archimedean energy correction?  
4. Finally, is there a stability property (intermediate between K‑polystability and uniform K‑stability) characterized by the sign, vanishing, or boundedness of \(C(\mathcal{X},\mathcal{L})\)?  

Why is it a "good" problem:  
This question isolates a precise and currently open refinement of the analytic–non‑Archimedean correspondence central to the Yau–Tian–Donaldson picture.  The slope equality capturing the leading‑order term is known, but the existence and meaning of a *constant‑term invariant* \(C(\mathcal{X},\mathcal{L})\) is completely unexplored.  Establishing a systematic expansion \(E^{\mathrm{an}}(\varphi_t)=tE^{\mathrm{NA}}+C+o(1)\) would extend the differential–algebraic bridge one order deeper, potentially revealing a new algebraic quantity governing "second‑order stability."  Such a discovery would refine the current hierarchy of stability notions and clarify the analytic asymptotics of energy functionals along degenerations.  Conceptually simple yet demanding deep analytic, pluripotential, and intersection‑theoretic insight, this problem embodies the “narrow entrance, profound depth’’ ideal of a good research‑level challenge.